\documentclass[12pt]{article}
\usepackage[utf8]{inputenc}
\usepackage[T1]{fontenc}
\usepackage[margin=1in]{geometry}
\usepackage{setspace}
\usepackage{xcolor}
\usepackage[colorlinks=true, urlcolor=blue!60!black, linkcolor=blue!60!black]{hyperref}
\setlength{\parindent}{2em}
\setlength{\parskip}{1.25em}
\renewcommand{\baselinestretch}{1.0}
\begin{document}
\begin{center}
{\LARGE \textit{Pierce v. Society of Sisters}, 268 U.S. 510 (1925)\par}
\vspace{0.8em}
{\large Case Brief --- Four-Page Edition\par}
\end{center}
\vspace{1em}

\section{Memory Jogger}

A unanimous Court struck down Oregon's compulsory public-school law under the Due Process Clause: the state cannot ``standardize its children'' by abolishing private and parochial schools, because parents hold a fundamental liberty interest in directing their children's upbringing and education (\href{https://supreme.justia.com/cases/federal/us/268/510/}{Justia}).

\section{Detailed Case Facts}

In 1922, amid post-war nativism and anti-Catholic sentiment backed by groups including the Ku Klux Klan, Oregon voters adopted the Compulsory Education Act by initiative. Effective in 1926, it required every parent of a child aged eight to sixteen to send the child to a public school, on pain of fines and jail, with narrow exemptions and none for private schools---effectively outlawing private and parochial primary education. Two corporations sued Governor Pierce and other officials: the Society of the Sisters of the Holy Names of Jesus and Mary, which ran Catholic schools and orphanages teaching standard subjects plus religious instruction, and Hill Military Academy, a for-profit secular school for boys. Officials' vows of strict enforcement were already causing parents to withdraw students and eroding the schools' income and property values.

\section{Procedural History}

The schools sought injunctions in the District of Oregon; the defendants did not contest the facts. A three-judge court granted preliminary injunctions, holding that the Fourteenth Amendment protected the schools' property and goodwill and that the suits were not premature. Oregon appealed directly, and the Supreme Court affirmed on June 1, 1925.

\section{Judicial Votes}

Affirmed, \textbf{9--0}.
\begin{itemize}
\item \textbf{Opinion of the Court:} Justice McReynolds, joined by Chief Justice Taft and Justices Holmes, Van Devanter, Brandeis, Sutherland, Butler, Sanford, and Stone.
\item \textbf{Concurrences or dissents:} None.
\end{itemize}

\section{Holding}

Under the Due Process Clause and \textit{Meyer v. Nebraska}, a state may not compel all normal children to attend public schools only, because doing so unreasonably interferes with parents' liberty to direct their children's education without a reasonable relation to any purpose within the state's competency. Although corporations cannot claim Fourteenth Amendment liberty, private schools have standing to seek relief before a statute takes effect when unconstitutional compulsion of their patrons threatens immediate, irreparable destruction of their business and property. The Court's concession that states may reasonably regulate all schools is technically dicta, since Oregon had abolished private schools rather than regulating them, but it later shaped \textit{Wisconsin v. Yoder}.

\section{Analysis of Opinions}

\textbf{Opinion of the Court (McReynolds).} The opinion blended \textit{Lochner}-era protection of property with an expansive vision of personal liberty. McReynolds first set aside what was not at issue: the state's power to require attendance at \emph{some} school and to inspect and set standards for all schools. Enforcing the Act would destroy the appellees' primary schools, a useful and meritorious enterprise, with no emergency justifying it. Relying on \textit{Meyer}, he held that the law interfered with parental liberty and bore no reasonable relation to a legitimate state purpose; the state may not monopolize instruction, for the child is not the State's creature. He then rejected three objections: the corporations had ``business and property'' threatened by the compulsion of their patrons, entitling them to equitable relief under \textit{Truax v. Raich} and \textit{Terrace v. Thompson}; the general law could not be defended as an exercise of the power to amend corporate charters (distinguishing \textit{Berea College}); and the suits were ripe because the threatened enforcement was already causing real and irreparable loss.

\section{Examples: Future Applications}

\begin{enumerate}
\item \textbf{Same-Side (Invalid): Public Middle Schools Only.} A state bans private middle schools to promote civic cohesion. Like Oregon's Act, it standardizes children through a state monopoly.
\item \textbf{Same-Side (Invalid): Permission to Board Elsewhere.} Parents may enroll children in out-of-state academies only if the local superintendent certifies the public school cannot meet the child's needs. Conditioning lawful private education on state permission forces families into public schools.
\item \textbf{Opposite-Side (Valid): Baseline Regulation.} Fire codes, immunization records, 180 instructional days, and core subjects including civics fall within the state's power to regulate all schools.
\item \textbf{Opposite-Side (Valid): No Subsidy.} Denying direct operating grants to religious schools (absent discriminatory exclusion from a general benefit) is valid; \textit{Pierce} protects a negative liberty, not a right to subsidy.
\item \textbf{Fence-Sitter: Substantial Equivalency.} A state requires five hours a day of secular instruction in English by certified teachers, threatening yeshivas focused on Talmudic study. The case sits between parental liberty, reinforced by \textit{Yoder}, and the regulatory dicta.
\end{enumerate}

\section{Critique}

\textbf{Originalist views.} \textit{Pierce} rests on an unenumerated liberty born in the \textit{Lochner} era. In \textit{Troxel}, Justice Scalia acknowledged parental authority as an unalienable right but called \textit{Meyer} and \textit{Pierce} relics of ``an era rich in substantive due process holdings that have since been repudiated.'' Others argue the case would rest more securely on the First Amendment, incorporated beginning one week later in \textit{Gitlow}.

\textbf{Progressive and child-rights views.} Progressives celebrate \textit{Pierce} as a cornerstone of pluralism and a seedbed for \textit{Griswold}, \textit{Roe}, \textit{Lawrence}, and \textit{Obergefell}. Barbara Bennett Woodhouse argued that \textit{Meyer} and \textit{Pierce} embedded a patriarchal view of the child as property, and James Dwyer questions whether robust parental control protects children's own interests. The case balances parental pluralism, democratic civic education, and children's independent interests.

\section{Key Quotations}

\begin{quote}
``Under the doctrine of \textit{Meyer v. Nebraska}, 262 U.\,S. 390, we think it entirely plain that the Act of 1922 unreasonably interferes with the liberty of parents and guardians to direct the upbringing and education of children under their control. As often heretofore pointed out, rights guaranteed by the Constitution may not be abridged by legislation which has no reasonable relation to some purpose within the competency of the State.''
\end{quote}
\noindent --- Justice McReynolds, for the Court, 268 U.S. at 534--35.

\begin{quote}
``The fundamental theory of liberty upon which all governments in this Union repose excludes any general power of the State to standardize its children by forcing them to accept instruction from public teachers only. The child is not the mere creature of the State; those who nurture him and direct his destiny have the right, coupled with the high duty, to recognize and prepare him for additional obligations.''
\end{quote}
\noindent --- Justice McReynolds, for the Court, 268 U.S. at 535.

\begin{quote}
``No question is raised concerning the power of the State reasonably to regulate all schools, to inspect, supervise and examine them, their teachers and pupils; to require that all children of proper age attend some school, that teachers shall be of good moral character and patriotic disposition, that certain studies plainly essential to good citizenship must be taught, and that nothing be taught which is manifestly inimical to the public welfare.''
\end{quote}
\noindent --- Justice McReynolds, for the Court, 268 U.S. at 534.

\end{document}
